\documentclass[5p]{elsarticle}
\usepackage{amsmath}
\usepackage{booktabs}
\usepackage{graphicx}
\usepackage{xurl}
\journal{Economics Letters}

\begin{document}

\begin{frontmatter}

\title{Border Premiums, License Caps, and Captive Demand in U.S.\ Cannabis Retail}

\author[1]{Mahyar Rezazadeh}
\author[1]{Babak Sahragard}
\address[1]{Department of Economics}

\begin{abstract}
We ask three questions about cannabis retail near U.S.\ state borders. First, does a border premium exist:
counties in Colorado and Washington within 70 miles of a nonlegal-state border earn 41 percent higher revenue
per dispensary than interior counties (county-clustered $p=0.033$), using both a binary threshold and a
continuous distance measure, with no comparable pattern in revenue per capita or a non-cannabis placebo
sector. Second, is this evidence of license-cap policy: Montana, the one sample state without a statewide
cap, shows the reverse pattern, a significant per-capita effect and a precise per-dispensary null. Third, how
is demand captured: the premium scales with the size of the captive nonlegal-side population within reach
($p=0.036$), not merely with proximity to it. Together these results identify a border premium, tie it to a
specific state policy mechanism, and show it responds to market size.
\end{abstract}

\begin{keyword}
cannabis \sep cross-border shopping \sep market entry \sep license caps \sep border economics
\end{keyword}

\end{frontmatter}

\noindent\textbf{Highlights}
\begin{itemize}
\item Revenue per dispensary is 41\% higher in U.S.\ counties near a nonlegal-state border ($p=0.033$).
\item A continuous log-distance measure gives the same result ($p=0.023$), not only the 70-mile cutoff.
\item The premium appears in revenue per store, not revenue per capita, where entry is capped (CO, WA).
\item Montana, with no statewide cap, shows the opposite pattern: significant per-capita, null per-store.
\item The premium rises with the captive nonlegal population within reach ($p=0.036$).
\end{itemize}

\section{Introduction}

Recreational cannabis legalization has proceeded unevenly across U.S.\ states, leaving many legal markets
directly adjacent to states where recreational sales remain prohibited. \citet{HansenMillerWeber2020} show
that Washington retailers near the Oregon border lost a substantial share of sales once Oregon's own market
opened, direct evidence that consumers cross state lines to buy cannabis. This is one instance of a general
pattern: cross-border shopping responds to price and legal differences, weakens with distance
\citep{AsplundFribergWilander2007,Lovenheim2008,BeattyLarsenSommervoll2009,ChandraHeadTappata2014}, and
responds to the size of the population that can reach the more permissive market
\citep{KnightSchiff2012,Agrawal2015}.

We ask three questions. \textbf{Result 1}: does a border premium exist, using all border counties in
Colorado and Washington as treatment and all interior counties as control, with both a binary and a
continuous distance measure? \textbf{Result 2}: is this evidence of license-cap policy, testable directly
because Montana, unlike Colorado and Washington, imposes no statewide numerical cap on dispensary licenses?
\textbf{Result 3}: how is this demand captured -- does the premium scale with the size of the captive
population on the nonlegal side, not just proximity to it? We answer yes to all three, which we read as
evidence for a specific, testable cross-border captive-demand mechanism rather than an incidental
correlation with border proximity.

\section{Data and Empirical Strategy}

We use a county-year panel for Colorado (2014--2025) and Washington (2015--2025), $n=787$ (42 of Colorado's
64 counties, 38 of Washington's 39, the rest never having an open dispensary), built from the Colorado
Department of Revenue's county sales series \citep{COMarijuanaSalesByCounty} and the Washington Liquor and
Cannabis Board's fiscal-year county sales reports \citep{WALCBSalesByCounty}. Retail sales began January 1,
2014 in Colorado \citep{COMedLicensees} and July 8, 2014 in Washington \citep{WSLCB2014FirstLicenses}.
Dispensary counts use each state's own current licensee registry, which reports county directly
\citep{COMedLicensees,WALCBLicenseeList,WALCBCannabisRetailers2026}; population is the Census Bureau's 2023
county estimate \citep{CensusPopEst2023}, held fixed across the panel. A county is a border county if its
Census Bureau internal point (not a population-weighted centroid) lies within 70 miles of the nearest county
in a state without legal recreational retail; 6 of 42 Colorado counties and 6 of 38 Washington counties meet
this definition, and the classification is fixed per county rather than updated as neighboring states
legalize during the sample. The continuous measure used in Result 1 is log(distance in miles) to that
nearest nonlegal county. $P_i$, used in Result 3, is the combined population of \emph{all} nonlegal-state
counties within 70 miles of county $i$, using the same fixed 2023 population vintage. We estimate
{\small
\begin{equation}
\ln\left(\text{Rev}_{ist}/N_{ist}\right) = \alpha + \beta_1\,\text{Border}_i + \beta_2\ln(P_i) + \gamma\ln(\text{Pop}_{ist}) + \theta_t + \delta_s + \varepsilon_{ist},
\end{equation}}
with year and state fixed effects and county-clustered standard errors. Montana's contrast case
(2022--2025, averaging 9.8 border and 19.5 interior counties per year, near Idaho, North Dakota, or Wyoming)
uses its Department of Revenue's county sales reports \citep{MTRevenueCountySales} and licensed-dispensary
list \citep{MTDispensaryLocations}; its adult-use market began January 1, 2022 under House Bill 701
\citep{MTHB701,MTRevenueCannabis}.

\section{Result 1: Does a Border Premium Exist?}

\begin{table*}[h]
\centering
\caption{Result 1: The Border Premium in Revenue per Dispensary}
\label{tab:result1}
\small
\begin{tabular}{lcccc}
\toprule
Specification & Coef. & SE (HC1) & SE (cl.) & $p$ (cl.) \\
\midrule
\multicolumn{5}{l}{\textit{Binary: border ($\le$70mi) vs.\ interior}} \\
Full controls$^{a}$ & $+0.341$ & $0.075$ & $0.161$ & $0.033$ \\
\quad excl.\ largest outlier county & $+0.302$ & $0.068$ & $0.147$ & $0.040$ \\
Colorado only / Washington only & $+0.552$ / $+0.233$ & -- & $0.199$ / $0.243$ & $0.005$ / $0.334$ \\
Border $\times$ log(revenue per capita) & $+0.046$ & -- & $0.219$ & $0.833$ \\
Placebo: general retail per capita (CO)$^{b}$ & $-0.006$ & -- & $0.089$ & $0.947$ \\
\midrule
\multicolumn{5}{l}{\textit{Continuous: log(distance to nonlegal border, miles)}} \\
Full controls$^{a}$ & $-0.215$ & $0.048$ & $0.095$ & $0.023$ \\
\bottomrule
\end{tabular}
\begin{minipage}{0.95\textwidth}
\vspace{0.2em}
\footnotesize $^{a}$Log(population), year and state fixed effects. $^{b}$Colorado Department of Revenue
county-industry filings \citep{CDORRetailByCounty}, $n=110$. Main sample $n=787$ (80 counties: 42 Colorado,
38 Washington; 12 border, 68 interior).
\end{minipage}
\end{table*}

Table~\ref{tab:result1} shows counties near a nonlegal border earn significantly more revenue per
dispensary, both with a binary 70-mile threshold ($p=0.033$) and a continuous distance measure ($p=0.023$);
the result is not driven by a single county and holds within Colorado alone. Revenue per capita shows no
comparable pattern, nor does a non-cannabis placebo sector, arguing against a generic border-commerce
explanation.

\section{Result 2: Evidence of License-Cap Policy}

The pattern above is exactly what a supply-constrained-entry story predicts: if the number of dispensaries
cannot adjust to demand, additional demand shows up as revenue per existing store rather than more stores or
higher revenue per capita \citep{BresnahanReiss1991,HollenbeckUetake2021}. Washington caps licenses statewide
at 334, allocated by lottery \citep{WSLCB2014LotteryPR}; Colorado has no statewide cap but gives local
governments authority under Amendment 64 to ban or limit retail marijuana, an option most municipalities have
exercised.

\begin{table*}[h]
\centering
\caption{Result 2: License-Cap Policy, Capped vs.\ Uncapped States}
\label{tab:result2}
\small
\begin{tabular}{lcc}
\toprule
Outcome & Coefficient (border effect) & $p$-value (cl.) \\
\midrule
\multicolumn{3}{l}{\textit{Colorado + Washington (license caps bind)}} \\
log(revenue per dispensary) & $+0.341$ & $0.033$ \\
log(revenue per capita) & $+0.046$ & $0.833$ \\
\midrule
\multicolumn{3}{l}{\textit{Montana (no statewide numerical cap)}} \\
log(revenue per capita) & $+0.654$ & $0.036$ \\
log(revenue per dispensary) & $+0.012$ & $0.958^{a}$ \\
\bottomrule
\end{tabular}
\begin{minipage}{0.95\textwidth}
\vspace{0.2em}
\footnotesize $^{a}$HC1 SE; too few Montana border counties ($n=23$) for county-clustered inference.
\end{minipage}
\end{table*}

Montana, the only sample state without a statewide numerical cap, offers a direct within-sample test.
Table~\ref{tab:result2} shows the predicted reversal: a significant per-capita border effect ($p=0.036$) and
a precise per-dispensary null ($p=0.958$). A single coefficient could be coincidental; a predicted reversal
in exactly the state whose institutional structure differs in the relevant way is harder to dismiss as one.

\section{Result 3: How Is Demand Captured?}

Results 1 and 2 do not by themselves show the premium reflects captive market size rather than an unrelated
correlate of border proximity. We test this using $P_i$, the population of \emph{all} nonlegal-state counties
within 70 miles of county $i$, not only the nearest one.

\begin{table}[h]
\centering
\caption{Result 3: Scaling with the Captive Nonlegal Population}
\label{tab:result3}
\small
\begin{tabular}{lcc}
\toprule
Specification & Coef. & $p$ (cl.) \\
\midrule
log($P_i$), full controls & $+0.030$ & $0.036$ \\
\quad excl.\ largest outlier county & $+0.024$ & $0.078$ \\
\bottomrule
\end{tabular}
\end{table}

Revenue per dispensary rises with $P_i$ ($p=0.036$; Table~\ref{tab:result3}), and this survives excluding the
largest outlier county. Combined with the continuous-distance result in Table~\ref{tab:result1}, this
indicates the premium responds continuously to both proximity to a nonlegal border and the size of the
population on the other side, not only to a binary border classification -- the signature of a genuine
captive cross-border demand channel rather than an unrelated border correlate.

\section{Conclusion}

Cannabis retailers near nonlegal-state borders earn more revenue per dispensary; where this shows up depends
systematically on whether entry is capped, and the effect scales with the size of the captive population
across the border. This adds a testable margin to the cross-border literature
\citep{HansenMillerWeber2020,HaoCowan2020}: state licensing policy and captive market size, not just distance,
determine whether cross-border cannabis demand appears as higher revenue at existing stores, as more stores,
or is left partly unmet. A supplementary entry-density test found no significant border gap in dispensaries
per capita once Washington's licensee counts are corrected from an earlier, incomplete registry; we regard
Result 2's direct per-capita/per-dispensary contrast, not entry density, as our evidence for the license-cap
mechanism.

\bibliographystyle{elsarticle-harv}
\bibliography{references}

\end{document}
